% Generated draft: requires mathematical and source review.
\documentclass[11pt]{article}
\usepackage[margin=25mm]{geometry}
\usepackage{fontspec}
\setmainfont{Latin Modern Roman}
\usepackage{amsmath,amssymb,mathtools,unicode-math}
\setmathfont{Latin Modern Math}
\usepackage{enumitem,needspace,xurl,hyperref}
\setlength{\parindent}{0pt}
\setlength{\parskip}{5pt}
\setlength{\emergencystretch}{3em}
\newcommand{\N}{\mathbb{N}}
\newcommand{\Z}{\mathbb{Z}}
\newcommand{\Q}{\mathbb{Q}}
\newcommand{\R}{\mathbb{R}}
\newcommand{\C}{\mathbb{C}}
\newcommand{\F}{\mathbb{F}}
\newcommand{\A}{\mathbb{A}}
\newcommand{\PP}{\mathbb{P}}
\newcommand{\E}{\mathbb{E}}
\newcommand{\eps}{\varepsilon}
\newcommand{\dd}{\,\mathrm{d}}
\begin{document}
\section*{UnsolvedMath ID: 6200104}
\textbf{Source problem code:} catalogue entry
\subsection{Definitions and mathematical statement}
Let \(Y\) be a compact connected finite-dimensional locally CAT\((0)\) space, with its intended intrinsic metric, and let \(X=\widetilde Y\) carry the lifted metric. The usual complete length-space hypotheses make \(X\) CAT\((0)\). The source assumes rank \(n\ge2\), but does not define rank. Maximum dimension of one embedded flat and a higher-rank condition on every complete geodesic are different notions and must not be silently identified.

Subject to the intended rank definition, the conjectured conclusion is the alternative
\[
X\cong X_1\times X_2\quad\text{nontrivially},
\qquad\text{or}\qquad X\text{ is isometric to a Euclidean building}.
\]
For metric factors, the product means
\[
d((x_1,x_2),(y_1,y_2))^2=d_1(x_1,y_1)^2+d_2(x_2,y_2)^2,
\]
with neither factor a single point. This is the metric interpretation of the source's Riemannian direct product even for nonsmooth complexes. A Euclidean building is meant with its standard apartment/Coxeter metric structure; no thickness or irreducibility restriction is added.

The source says the problem is most natural for piecewise-Euclidean metric cell complexes; it does not explicitly limit the initial class to those complexes. The rank convention, any geodesic-extension/completeness condition, and whether that context is an actual hypothesis must be checked. In particular, the broad wording has no symmetric-space alternative, a distinction that matters when comparing it with general rank-rigidity formulations.

The source reports the conjecture proved by Ballmann--Brin in the cases of two- and three-dimensional complexes, citing [2,3]. These are dimensions of the complexes, not a claim about every compact space of rank two or three. The supplied research summary reports no verified all-dimensional resolution.

A trailing paragraph in the same extracted record begins a different topic, cogrowth. It defines cogrowth of \(H\le G\), for finitely generated \(G\), as growth of the Schreier graph associated to \(G/H\) and a Cayley graph of \(G\). This paragraph supplies no question and no hypothesis for the rank-rigidity assertion; it is preserved here as adjacent source context.
\subsection{Short English statement}
Does every intended higher-rank compact locally CAT\((0)\) space have universal cover splitting as a nontrivial metric product or forming a Euclidean building? The rank and space-class assumptions need confirmation; the source reports the two- and three-dimensional complex cases proved.
\subsection{Sources}
\begin{itemize}
\item Local UnsolvedMath registry, permanent ID \(6200104\), exact source code \texttt{catalogue entry}, ``Boundaries of Groups and Kleinian Groups --- Problem 104.'' Proposers Werner Ballmann and Misha Brin, 2005.
\item Kapovich, \emph{Problems on Boundaries of Groups and Kleinian Groups} (2005), Problem 104, recorded PDF page 26; \url{https://www.math.ucdavis.edu/~kapovich/EPR/problems.pdf}. Recorded access 2026-07-29; pdf-text-normalized extraction.
\item Source references [2,3], attributed to Ballmann--Brin for the two- and three-dimensional complex cases, lack their bibliography entries in this record. No complete later resolution reference is supplied.
\end{itemize}
\end{document}
